\chapter{The Sherrington--Kirkpatrick model and mixed $p$-spin models}\label{chap:sg}

This chapter formalizes Section~1.1 of Panchenko, \emph{The Sherrington--Kirkpatrick Model} (2013): the random
Hamiltonians of the Sherrington--Kirkpatrick (SK) model and of the pure and mixed $p$-spin models, their
covariance, the free energy and its smooth approximation of the ground-state energy, the monotonicity of the
free energy in the temperature, the zero-temperature limit, the Gibbs measure and the temperature derivative of
the free energy. The existence of the thermodynamic limit of the free energy (the Guerra--Toninelli theorem)
is \emph{not} formalized; it enters the theorems below only as a hypothesis. The Lean files are in
\texttt{SpinGlasses/} (modules \texttt{Basic}, \texttt{SKModel}, \texttt{FreeEnergy}, \texttt{ThermodynamicLimit},
\texttt{GibbsMeasure}, \texttt{MixedPSpin}), together with the general tools
\texttt{ForMathlib/Probability/Moments/GaussianPiLinearForm.lean}, \texttt{GaussianSeries.lean},
\texttt{ForMathlib/Probability/Independence/IntegralTsumMul.lean},
\texttt{ForMathlib/Topology/UniformSpace/LimitInterchange.lean} and
\texttt{ForMathlib/MeasureTheory/MeasurableSpace/Units.lean}.

\medskip\noindent\textbf{Notation.}
$\iota$ is a finite index set and $N=|\iota|$ (the number of spins; $\iota=\mathrm{Fin}\,N$ in the SK model, but all
statements hold for any finite type). A \emph{configuration} is $\sigma\in\Sigma_\iota=\{-1,+1\}^\iota$, a function
$\sigma:\iota\to\mathbb Z^\times$ in Lean (here $\mathbb Z^\times=\{\pm1\}$), read as the real vector
$(\sigma_i)_{i\in\iota}$ when sums are formed. The \emph{overlap} of two configurations is
$R(\sigma^1,\sigma^2)=N^{-1}\sum_i\sigma^1_i\sigma^2_i=N^{-1}\,\sigma^1\!\cdot\!\sigma^2$. The disorder
$G=(g_{ij})$ has i.i.d.\ $\mathcal N(0,1)$ entries, the law
\texttt{Measure.pi fun \_ => Measure.pi fun \_ => gaussianReal 0 1}, and the SK Hamiltonian is
$H_N(\sigma)=N^{-1/2}\sum_{i,j}g_{ij}\sigma_i\sigma_j=N^{-1/2}\,\sigma^\top G\sigma$. For an energy
$H:\Sigma\to\mathbb R$ and $\beta>0$ the partition function is $Z(\beta)=\sum_\sigma e^{\beta H(\sigma)}$ and the free
energy is $F_N(\beta)=N^{-1}\mathbb E\log Z_N(\beta)$.

\medskip\noindent\textbf{Which definitions are not introduced.}
Following the repository's convention this chapter introduces \emph{no} new Lean definition. The overlap, the
Hamiltonians, the partition function, the free energy, the covariance function $\xi$, the Gibbs weights and the
Gibbs measure are all written out explicitly in the statements: the overlap is
$N^{-1}\,\sigma\cdot\tau$ (an expression in \texttt{dotProduct}), the partition function is the sum
$\sum_\sigma e^{\beta H(\sigma)}$, the maximum is the repository's \texttt{Function.max}, and the Gibbs measure is
Mathlib's exponentially tilted counting measure \texttt{Measure.count.tilted} with tilt $\beta H$. Hamiltonians appear as an arbitrary function $H$ together with its
defining equation as a hypothesis (written \texttt{hH}).

\section{Configurations and overlap}\label{sec:sg:basic}

\begin{lemma}[Overlap]\label{lem:sg:overlap}
  \lean{SpinGlasses.card_inv_mul_dotProduct_self_eq_one, SpinGlasses.abs_card_inv_mul_dotProduct_le_one,
        SpinGlasses.card_spin_configurations, Matrix.mulVec_dotProduct_mulVec_of_transpose_mul_self}
  \leanok
Let $\iota$ be a nonempty finite type. (1) If $\sigma_i^2=1$ for all $i$ then $R(\sigma,\sigma)=1$.
(2) If $|\sigma_i|\le1$ and $|\tau_i|\le1$ for all $i$ then $|R(\sigma,\tau)|\le1$.
(3) $|\Sigma_\iota|=2^N$. (4) If $A^\top A=1$ for a matrix $A$ then $(Ax)\cdot(Ay)=x\cdot y$; in particular
$R(O\sigma,O\tau)=R(\sigma,\tau)$ for every orthogonal $O$.
\end{lemma}

\begin{proof}
  \leanok
(1) $\sum_i\sigma_i^2=N$. (2) $|\sum_i\sigma_i\tau_i|\le\sum_i|\sigma_i||\tau_i|\le N$ termwise, so no Cauchy--Schwarz
is needed. (3) $|\{\pm1\}^\iota|=2^{|\iota|}$ (\texttt{Fintype.card\_fun}). (4) $(Ax)\cdot(Ay)=x\cdot(A^\top A)y=x\cdot y$.
\end{proof}

\section{The SK Hamiltonian and its covariance}\label{sec:sg:sk}

\begin{lemma}[Covariance of linear forms of a Gaussian matrix]\label{lem:sg:bilinear}
  \lean{NTK.dotProduct_mulVec_of_eq_sum, NTK.memLp_linearForm_gaussianInit,
        NTK.integral_linearForm_gaussianInit, NTK.integral_gaussianMatrix_bilinear_mul_bilinear}
  \leanok
Let $G\in\mathbb R^{\iota\times\kappa}$ have i.i.d.\ $\mathcal N(0,1)$ entries and let $u,x\in\mathbb R^\iota$,
$v,y\in\mathbb R^\kappa$ be deterministic. Then $u^\top Gv=\sum_{i,j}u_iv_jg_{ij}$ is square-integrable and centered, and
\[\mathbb E\big[(u^\top Gv)(x^\top Gy)\big]=(u\cdot x)(v\cdot y).\]
\end{lemma}

\begin{proof}
  \leanok
Write $G$ as the random vector $q\mapsto g_q$ indexed by $q=(i,j)\in\iota\times\kappa$. Its coordinates are
independent standard Gaussians, so the second moment matrix is the identity
(\texttt{integral\_gaussianInit\_entry\_mul\_entry}), and for any random vector $X$ with square-integrable
coordinates $\mathbb E[(X\cdot a)(X\cdot b)]=a^\top\Gamma b$ with $\Gamma=\mathbb E[XX^\top]$
(\texttt{ProbabilityTheory.integral\_dotProduct\_mul\_dotProduct}). With $a_{(i,j)}=u_iv_j$ and
$b_{(i,j)}=x_iy_j$ we get $\sum_{i,j}u_iv_jx_iy_j=(u\cdot x)(v\cdot y)$.
\end{proof}

\begin{theorem}[Covariance of the SK Hamiltonian, Panchenko Eq.~(1.3)]\label{thm:sg:skcov}
  \lean{SpinGlasses.integral_skHamiltonian_mul_skHamiltonian}
  \leanok
For all real vectors $\sigma,\tau\in\mathbb R^\iota$,
$\mathbb E[H_N(\sigma)H_N(\tau)]=N\,R(\sigma,\tau)^2$.
\end{theorem}

\begin{proof}
  \leanok
  \uses{lem:sg:bilinear}
Apply Lemma~\ref{lem:sg:bilinear} with $u=v=\sigma$ and $x=y=\tau$: the covariance of the unnormalized forms is
$(\sigma\cdot\tau)^2$. The normalization $N^{-1/2}$ appears twice, giving
$N^{-1}(\sigma\cdot\tau)^2=N\,(N^{-1}\sigma\cdot\tau)^2$.
\end{proof}

\begin{lemma}[Dean's comfort identity]\label{lem:sg:comfort}
  \lean{SpinGlasses.sum_lt_mul_units_eq_sub}
  \leanok
For a totally ordered index set and $\sigma\in\{\pm1\}^\iota$, with the nested sums over $i<j$,
\[\sum_{i<j}g_{ij}\sigma_i\sigma_j=\sum_{i<j,\ \sigma_i=\sigma_j}g_{ij}-\sum_{i<j,\ \sigma_i\ne\sigma_j}g_{ij}.\]
\end{lemma}

\begin{proof}
  \leanok
For $\sigma_i,\sigma_j\in\{\pm1\}$ the product $\sigma_i\sigma_j$ equals $+1$ if $\sigma_i=\sigma_j$ and $-1$
otherwise; split each term accordingly.
\end{proof}

\section{Free energy and the smooth approximation of the ground-state energy}\label{sec:sg:free}

The estimate of this section is deterministic; it is proved first for a finite set, then integrated, then specialized
to the hypercube, so that it applies to any finite Gibbs system and not only to SK.

\begin{lemma}[Log-sum-exp sandwich]\label{lem:sg:lse}
  \lean{SpinGlasses.max_le_inv_mul_log_sum_exp, SpinGlasses.inv_mul_log_sum_exp_le,
        SpinGlasses.abs_inv_mul_log_sum_exp_sub_max_le}
  \leanok
For a finite nonempty set $\alpha$, $f:\alpha\to\mathbb R$ and $\beta>0$,
\[\max_a f(a)\le\frac1\beta\log\sum_a e^{\beta f(a)}\le\frac{\log|\alpha|}{\beta}+\max_a f(a).\]
\end{lemma}

\begin{proof}
  \leanok
The sum contains the term $e^{\beta\max f}$, and it is at most $|\alpha|e^{\beta\max f}$. Take logarithms and divide
by $\beta>0$.
\end{proof}

\begin{lemma}[Integrability of the random landscape]\label{lem:sg:integrable}
  \lean{SpinGlasses.integrable_max, SpinGlasses.measurable_log_sum_exp, SpinGlasses.integrable_log_sum_exp}
  \leanok
Let $H:\Omega\to\mathbb R^\alpha$ have measurable, integrable coordinates $\omega\mapsto H(\omega)_a$ on a finite
measure space, $\alpha$ finite nonempty and $\beta>0$. Then $\omega\mapsto\max_aH(\omega)_a$ and
$\omega\mapsto\log\sum_ae^{\beta H(\omega)_a}$ are integrable.
\end{lemma}

\begin{proof}
  \leanok
  \uses{lem:sg:lse}
$|\max_aH_a|\le\sum_a|H_a|$, which is integrable. By Lemma~\ref{lem:sg:lse},
$|\log Z|\le\beta|\max H|+\log|\alpha|$, again integrable on a finite measure space; measurability is that of a
finite composition.
\end{proof}

\begin{theorem}[Smooth approximation, Panchenko Eq.~(1.7)]\label{thm:sg:smooth}
  \lean{SpinGlasses.integral_max_le_inv_mul_integral_log_sum_exp,
        SpinGlasses.inv_mul_integral_log_sum_exp_le,
        SpinGlasses.abs_inv_mul_integral_log_sum_exp_sub_integral_max_le,
        SpinGlasses.log_card_spin_configurations, SpinGlasses.smooth_approximation_sandwich,
        SpinGlasses.smooth_approximation_error_bound}
  \leanok
Let $H$ be a random energy on $\Sigma_\iota$ with measurable integrable coordinates under a probability measure, and
$\beta>0$. Then
\[\frac1N\mathbb E\max_\sigma H(\sigma)\le\frac{1}{N\beta}\mathbb E\log Z(\beta)\le
\frac{\log2}{\beta}+\frac1N\mathbb E\max_\sigma H(\sigma),\]
hence $\big|\frac1{N\beta}\mathbb E\log Z(\beta)-\frac1N\mathbb E\max_\sigma H(\sigma)\big|\le\frac{\log2}\beta$.
(For $N=0$ both sides vanish by the convention $0^{-1}=0$.)
\end{theorem}

\begin{proof}
  \leanok
  \uses{lem:sg:lse, lem:sg:integrable}
Integrate Lemma~\ref{lem:sg:lse} (Bochner integrals are monotone, and the integrands are integrable by
Lemma~\ref{lem:sg:integrable}), then use $\log|\Sigma_\iota|=N\log2$ (Lemma~\ref{lem:sg:overlap}(3)) and divide
by $N$.
\end{proof}

\begin{theorem}[Smooth approximation for the SK model]\label{thm:sg:sksmooth}
  \lean{SpinGlasses.sk_measurable_integrable, SpinGlasses.sk_smooth_approximation_sandwich,
        SpinGlasses.sk_smooth_approximation_error_bound}
  \leanok
For the SK Hamiltonian under i.i.d.\ Gaussian disorder, the hypotheses of Theorem~\ref{thm:sg:smooth} hold, so its
conclusion holds with $H=H_N$.
\end{theorem}

\begin{proof}
  \leanok
  \uses{thm:sg:smooth, lem:sg:bilinear}
For fixed $\sigma$, $G\mapsto H_N(\sigma)$ is a finite linear combination of the coordinates $g_{ij}$, hence
measurable, and it is integrable because it is square-integrable (Lemma~\ref{lem:sg:bilinear}).
\end{proof}

\section{Monotonicity in the temperature and the zero-temperature limit}\label{sec:sg:thermo}

\begin{lemma}[Lyapunov monotonicity of the log-partition function]\label{lem:sg:lyapunov}
  \lean{SpinGlasses.monotoneOn_inv_mul_log_sum_exp_sub_log_card,
        SpinGlasses.monotoneOn_inv_mul_integral_log_sum_exp_sub_log_card,
        SpinGlasses.monotoneOn_inv_mul_free_energy_sub_log_two}
  \leanok
For a finite nonempty set $\alpha$ and $f:\alpha\to\mathbb R$, the function
$\beta\mapsto\beta^{-1}\big(\log\sum_ae^{\beta f(a)}-\log|\alpha|\big)$ is non-decreasing on $(0,\infty)$. The same
holds for $\beta\mapsto\beta^{-1}\big(\mathbb E\log\sum_ae^{\beta H_a}-\log|\alpha|\big)$ for a random energy as in
Lemma~\ref{lem:sg:integrable}, and for a random energy on the hypercube the function
$\beta\mapsto\beta^{-1}\big(F_N(\beta)-\log2\big)$ is non-decreasing on $(0,\infty)$.
\end{lemma}

\begin{proof}
  \leanok
  \uses{lem:sg:integrable}
Let $\mu_0$ be the uniform probability measure on $\alpha$. Then $\beta^{-1}(\log\sum_ae^{\beta f}-\log|\alpha|)=
\log\|e^f\|_{L^\beta(\mu_0)}$, and the claim is that the $L^p$ norm increases with $p$. Directly: for
$0<\beta_1\le\beta_2$ set $p=\beta_2/\beta_1\ge1$ and apply the power-mean inequality
$\big(\sum_a|\alpha|^{-1}z_a\big)^p\le\sum_a|\alpha|^{-1}z_a^p$ to $z_a=e^{\beta_1f(a)}$. Taking logarithms gives
$\frac{\beta_2}{\beta_1}\log\mathrm{mean}(e^{\beta_1f})\le\log\mathrm{mean}(e^{\beta_2f})$; divide by $\beta_2$. The
random version integrates the pointwise inequality, and the hypercube version divides by $N$ and uses
$\log|\Sigma_\iota|=N\log2$.
\end{proof}

\begin{theorem}[Thermodynamic limit and zero-temperature limit, Panchenko Eq.~(1.10)]\label{thm:sg:zerotemp}
  \lean{exists_tendsto_of_dist_le_of_tendsto, SpinGlasses.monotoneOn_inv_mul_limit_sub_log_two,
        SpinGlasses.exists_tendsto_ground_state_energy, SpinGlasses.sk_monotoneOn_inv_mul_limit_sub_log_two,
        SpinGlasses.sk_exists_tendsto_ground_state_energy}
  \leanok
Consider a family of hypercube models indexed by $k\to\infty$ (for SK, $k=N$ and $\iota=\mathrm{Fin}\,N$), with
measurable, integrable coordinates, and suppose that for every $\beta>0$ the free energy converges,
$F_k(\beta)\to F(\beta)$. Then:
(1) $\beta\mapsto\beta^{-1}(F(\beta)-\log2)$ is non-decreasing on $(0,\infty)$;
(2) the normalized ground-state energies $E_k=|\iota_k|^{-1}\mathbb E\max_\sigma H_k(\sigma)$ converge to some $E_0$, and
$\beta^{-1}F(\beta)\to E_0$ as $\beta\to\infty$.
\end{theorem}

\begin{proof}
  \leanok
  \uses{lem:sg:lyapunov, thm:sg:smooth}
(1) A pointwise limit of monotone functions is monotone (Mathlib's
\texttt{monotoneOn\_of\_frequently\_monotoneOn\_of\_tendsto}) applied to Lemma~\ref{lem:sg:lyapunov}.
(2) By Theorem~\ref{thm:sg:smooth}, $|\beta^{-1}F_k(\beta)-E_k|\le\varepsilon(\beta):=\log2/\beta$ uniformly in $k$. The
limit $g(\beta)=\lim_kF_k(\beta)/\beta$ then satisfies $|g(\beta)-g(\beta')|\le\varepsilon(\beta)+\varepsilon(\beta')$ (let
$k\to\infty$ in the triangle inequality), so $g$ is Cauchy as $\beta\to\infty$ and has a limit $E_0$; the uniform
approximation then makes $E_k$ converge to the same $E_0$ (interchange of limits, in the abstract form of
\texttt{exists\_tendsto\_of\_dist\_le\_of\_tendsto}). In particular, unlike the statement in Panchenko, the
convergence of the ground-state energy is a conclusion and not an assumption.
\end{proof}

\section{The Gibbs measure}\label{sec:sg:gibbs}

For $\beta\in\mathbb R$ and $H:\alpha\to\mathbb R$ on a finite set, the \emph{Gibbs measure}
$G_\beta(a)=e^{\beta H(a)}/Z(\beta)$ is the exponential tilt of the counting measure by $\beta H$; in Lean it is
\texttt{Measure.count.tilted} with tilt $\beta H$, and it is a probability measure whenever $\alpha$ is nonempty. We write
$\langle g\rangle_\beta=\int g\,dG_\beta$.

\begin{lemma}[Gibbs averages]\label{lem:sg:gibbsavg}
  \lean{SpinGlasses.integral_count_eq_sum, SpinGlasses.mgf_count, SpinGlasses.integral_count_tilted,
        SpinGlasses.real_count_tilted_singleton, SpinGlasses.norm_integral_count_tilted_le}
  \leanok
(1) $Z(\beta)=\sum_ae^{\beta H(a)}$ is the moment-generating function of $H$ for the counting measure.
(2) $\langle g\rangle_\beta=\sum_a\frac{e^{\beta H(a)}}{Z(\beta)}\,g(a)$ and $G_\beta(\{a\})=e^{\beta H(a)}/Z(\beta)$.
(3) $|\langle g\rangle_\beta|\le\sum_a|g(a)|$.
\end{lemma}

\begin{proof}
  \leanok
(1) The integral against counting measure on a finite set is the sum. (2) Mathlib's
\texttt{integral\_tilted\_mul\_eq\_mgf}, and the indicator of $\{a\}$ for the mass of the singleton.
(3) $G_\beta$ is a probability measure and $|g|\le\sum_a|g(a)|$ pointwise.
\end{proof}

\begin{theorem}[Temperature derivative of the free energy]\label{thm:sg:deriv}
  \lean{SpinGlasses.hasDerivAt_log_sum_exp, SpinGlasses.hasDerivAt_integral_log_sum_exp,
        SpinGlasses.sk_hasDerivAt_free_energy}
  \leanok
(1) For every $H:\alpha\to\mathbb R$, $\frac{d}{d\beta}\log Z(\beta)=\langle H\rangle_\beta$.
(2) For a random energy with measurable, integrable coordinates and $\beta>0$,
$\frac{d}{d\beta}\mathbb E\log Z(\beta)=\mathbb E\langle H\rangle_\beta$.
(3) For the SK model, $\frac d{d\beta}F_N(\beta)=N^{-1}\mathbb E\langle H_N\rangle_\beta$ for $\beta>0$.
\end{theorem}

\begin{proof}
  \leanok
  \uses{lem:sg:gibbsavg, lem:sg:integrable}
(1) $\log Z=\mathrm{cgf}_H$, the cumulant-generating function of $H$ for the counting measure; every exponential moment
is finite on a finite set, so Mathlib's \texttt{integral\_tilted\_mul\_self} gives
$\mathrm{cgf}_H'(\beta)=\int H\,d(\mu.\mathrm{tilted}(\beta H))$. (2) Differentiate under the integral sign
(dominated differentiation): the pointwise derivative is $\langle H(\omega)\rangle_\beta$ by (1), and
$|\langle H(\omega)\rangle_\beta|\le\sum_a|H(\omega)_a|$ uniformly in $\beta$ by Lemma~\ref{lem:sg:gibbsavg}(3), which
is integrable by assumption. (3) is (2) with the SK landscape (Theorem~\ref{thm:sg:sksmooth}), multiplied by $N^{-1}$.
\end{proof}

\begin{theorem}[Zero-temperature concentration]\label{thm:sg:concentration}
  \lean{SpinGlasses.tendsto_real_count_tilted_singleton_of_lt, SpinGlasses.tendsto_integral_count_tilted_max}
  \leanok
Let $\alpha$ be finite nonempty and $H:\alpha\to\mathbb R$. (1) If $H(a)<\max H$ then $G_\beta(\{a\})\to0$ as
$\beta\to\infty$; indeed $G_\beta(\{a\})\le e^{-\beta(\max H-H(a))}$. (2) $\langle H\rangle_\beta\to\max H$.
\end{theorem}

\begin{proof}
  \leanok
  \uses{lem:sg:gibbsavg}
(1) $Z(\beta)\ge e^{\beta\max H}$ since it contains the term of a maximizer, hence
$e^{\beta H(a)}/Z(\beta)\le e^{-\beta(\max H-H(a))}\to0$. (2) $\max H-\langle H\rangle_\beta=\sum_aG_\beta(\{a\})(\max H-H(a))$
because the weights sum to one; every term with $H(a)<\max H$ tends to $0$ by (1) and the others vanish identically.
\end{proof}

\begin{remark}
The convergence of the Gibbs weights of the ground states to $1/|\mathrm{argmax}\,H|$ is not formalized, since it
is not used anywhere in this chapter.
\end{remark}

\section{Pure and mixed $p$-spin models}\label{sec:sg:pspin}

For $p\ge1$ the pure $p$-spin Hamiltonian is
\[H_{N,p}(\sigma)=N^{-(p-1)/2}\sum_{i_1,\dots,i_p\in\iota}g_{i_1\dots i_p}\sigma_{i_1}\cdots\sigma_{i_p},\]
with i.i.d.\ standard Gaussian disorder $g:(\mathrm{Fin}\,p\to\iota)\to\mathbb R$ (the case $p=2$ is the SK
Hamiltonian). The mixed $p$-spin Hamiltonian is $H_N(\sigma)=\sum_{p\ge1}\beta_pH_{N,p}(\sigma)$ with independent
disorders $g_p$, i.e.\ $g\sim\bigotimes_p\mathcal N(0,1)^{\iota^p}$ (\texttt{Measure.infinitePi}), and the covariance
function is $\xi(x)=\sum_{p\ge1}\beta_p^2x^p$. In Lean the condition $\beta_0=0$ replaces ``the sum starts at $p=1$''.

\begin{lemma}[Gaussian linear forms and series]\label{lem:sg:gaussseries}
  \lean{ProbabilityTheory.integral_dotProduct_mul_dotProduct_pi_gaussianReal,
        ProbabilityTheory.integral_dotProduct_pi_gaussianReal,
        ProbabilityTheory.memLp_dotProduct_pi_gaussianReal,
        ProbabilityTheory.integral_tsum_mul_tsum_of_indepFun,
        ProbabilityTheory.integrable_tsum_of_tsum_eLpNorm_two_ne_top,
        ProbabilityTheory.eLpNorm_two_le_ofReal_of_integral_sq_le,
        ProbabilityTheory.integral_tsum_dotProduct_mul_tsum_dotProduct,
        ProbabilityTheory.measurable_tsum_dotProduct, ProbabilityTheory.integrable_tsum_dotProduct}
  \leanok
(1) For $g$ with i.i.d.\ $\mathcal N(0,1)$ coordinates indexed by a finite set $\kappa$ and $a,b\in\mathbb R^\kappa$, the
linear forms $a\cdot g$ are square-integrable, centered, and $\mathbb E[(a\cdot g)(b\cdot g)]=a\cdot b$.
(2) Let $X_p,Y_p$ ($p$ in a countable set) be square-integrable with $\sum_p\|X_p\|_{L^2}<\infty$,
$\sum_p\|Y_p\|_{L^2}<\infty$, each $X_p$ centered and independent of $Y_q$ for $p\ne q$. Then
$\sum_pX_p$ is integrable and $\mathbb E[(\sum_pX_p)(\sum_qY_q)]=\sum_p\mathbb E[X_pY_p]$.
(3) Let $g_p$ ($p$ countable) be independent with i.i.d.\ $\mathcal N(0,1)$ coordinates and let $a_p,b_p$ be coefficient
vectors with $\sum_p\|a_p\|_2<\infty$ and $\sum_p\|b_p\|_2<\infty$. Then $\sum_pa_p\cdot g_p$ is measurable and integrable and
$\mathbb E[(\sum_pa_p\cdot g_p)(\sum_pb_p\cdot g_p)]=\sum_pa_p\cdot b_p$.
\end{lemma}

\begin{proof}
  \leanok
(1) The law of $g$ is the standard Gaussian of $\mathbb R^\kappa$, whose covariance bilinear form is the inner product,
and $\mathbb E g=0$. (2) The series converge absolutely almost surely, since $\sum_p\|X_p\|_{L^1}\le\sum_p\|X_p\|_{L^2}$
(Mathlib's \texttt{summable\_norm\_of\_tsum\_eLpNorm\_ne\_top}), so $(\sum_pX_p)(\sum_qY_q)=\sum_{(p,q)}X_pY_q$ almost
everywhere. By H\"older, $\mathbb E|X_pY_q|\le\|X_p\|_{L^2}\|Y_q\|_{L^2}$, and $\sum_{p,q}\|X_p\|_{L^2}\|Y_q\|_{L^2}<\infty$, so the
double series may be integrated termwise (\texttt{integral\_tsum}). For $p\ne q$, independence and centering give
$\mathbb E[X_pY_q]=\mathbb E[X_p]\mathbb E[Y_q]=0$, which leaves the diagonal. (3) Apply (2) to $X_p=a_p\cdot g_p$ and
$Y_p=b_p\cdot g_p$: they are independent for $p\ne q$ because they are functions of different coordinates of the
infinite product, $\|X_p\|_{L^2}=\|a_p\|_2$ by (1), and $\mathbb E[X_pY_p]=a_p\cdot b_p$ by (1).
\end{proof}

\begin{theorem}[Covariance of the pure $p$-spin Hamiltonian, Panchenko Eq.~(1.13)]\label{thm:sg:purecov}
  \lean{SpinGlasses.sum_prod_mul_prod_eq_pow, SpinGlasses.pureCoeff_dotProduct_pureCoeff,
        SpinGlasses.integral_purePSpin_mul_purePSpin}
  \leanok
For $p\ge1$, a nonempty finite $\iota$ and any $\sigma,\tau\in\mathbb R^\iota$,
$\mathbb E[H_{N,p}(\sigma)H_{N,p}(\tau)]=N\,R(\sigma,\tau)^p$.
\end{theorem}

\begin{proof}
  \leanok
  \uses{lem:sg:gaussseries}
$H_{N,p}(\sigma)=u(\sigma)\cdot g$ with $u(\sigma)_{i_1\dots i_p}=N^{-(p-1)/2}\prod_k\sigma_{i_k}$. By
Lemma~\ref{lem:sg:gaussseries}(1) the covariance is $u(\sigma)\cdot u(\tau)=N^{-(p-1)}\sum_{i_1\dots i_p}\prod_k\sigma_{i_k}\tau_{i_k}
=N^{-(p-1)}(\sigma\cdot\tau)^p$, since the sum over multi-indices of a product factorizes
(\texttt{Fintype.sum\_pow}). This equals $N(N^{-1}\sigma\cdot\tau)^p$ for $p\ge1$.
\end{proof}

\begin{lemma}[Summability]\label{lem:sg:summable}
  \lean{SpinGlasses.summable_abs_of_summable_two_pow_mul_sq, SpinGlasses.summable_sq_of_summable_abs,
        SpinGlasses.summable_sq_mul_pow_of_summable_abs}
  \leanok
(1) If $\sum_p2^p\beta_p^2<\infty$ (Panchenko's condition) then $\sum_p|\beta_p|<\infty$.
(2) If $\sum_p|\beta_p|<\infty$ then $\sum_p\beta_p^2<\infty$ and $\xi(x)=\sum_p\beta_p^2x^p$ converges absolutely for
$|x|\le1$.
\end{lemma}

\begin{proof}
  \leanok
(1) $2|\beta_p|\le2^p\beta_p^2+2^{-p}$, since $(2^p|\beta_p|-1)^2\ge0$, and $\sum_p2^{-p}<\infty$. (2) $\beta_p^2\le
(\sum_q|\beta_q|)\,|\beta_p|$ and $|\beta_p^2x^p|\le\beta_p^2$.
\end{proof}

\begin{theorem}[Covariance of the mixed $p$-spin Hamiltonian, Panchenko Eq.~(1.14)]\label{thm:sg:mixedcov}
  \lean{SpinGlasses.mixedPSpin_measurable, SpinGlasses.mixedPSpin_integrable,
        SpinGlasses.integral_mixedPSpin_mul_mixedPSpin}
  \leanok
Let $\iota$ be nonempty and finite, $\beta_0=0$, $\sum_p|\beta_p|<\infty$ and $\sigma,\tau\in[-1,1]^\iota$. Then the
series $H_N(\sigma)=\sum_p\beta_pH_{N,p}(\sigma)$ is measurable and integrable in the disorder, and
\[\mathbb E[H_N(\sigma)H_N(\tau)]=N\,\xi\big(R(\sigma,\tau)\big).\]
\end{theorem}

\begin{proof}
  \leanok
  \uses{lem:sg:gaussseries, thm:sg:purecov, lem:sg:summable}
$H_N(\sigma)=\sum_pa_p(\sigma)\cdot g_p$ with $a_p(\sigma)=\beta_pu_p(\sigma)$. By Theorem~\ref{thm:sg:purecov},
$\|a_p(\sigma)\|_2^2=\beta_p^2NR(\sigma,\sigma)^p\le\beta_p^2N$ because $0\le R(\sigma,\sigma)\le1$ for
$\sigma\in[-1,1]^\iota$; so $\|a_p(\sigma)\|_2\le|\beta_p|\sqrt N$ is summable in $p$. Lemma~\ref{lem:sg:gaussseries}(3) then gives
measurability, integrability and $\mathbb E[H_N(\sigma)H_N(\tau)]=\sum_pa_p(\sigma)\cdot a_p(\tau)
=\sum_p\beta_p^2NR(\sigma,\tau)^p$ (for $p=0$ both sides vanish because $\beta_0=0$).
\end{proof}

\begin{theorem}[Free energy of the mixed $p$-spin model]\label{thm:sg:mixedfree}
  \lean{SpinGlasses.mixedPSpin_smooth_approximation_error_bound,
        SpinGlasses.mixedPSpin_hasDerivAt_free_energy}
  \leanok
Under the hypotheses of Theorem~\ref{thm:sg:mixedcov}, for $t>0$,
$\big|\frac1{Nt}\mathbb E\log Z_N(t)-\frac1N\mathbb E\max_\sigma H_N(\sigma)\big|\le\frac{\log2}{t}$ and
$\frac{d}{dt}\frac1N\mathbb E\log Z_N(t)=\frac1N\mathbb E\langle H_N\rangle_t$, where $Z_N(t)=\sum_\sigma e^{tH_N(\sigma)}$.
\end{theorem}

\begin{proof}
  \leanok
  \uses{thm:sg:mixedcov, thm:sg:smooth, thm:sg:deriv}
For a hypercube configuration the coordinates are in $[-1,1]$, so Theorem~\ref{thm:sg:mixedcov} supplies the
measurability and integrability of $G\mapsto H_N(\sigma)$ required by Theorems~\ref{thm:sg:smooth} and~\ref{thm:sg:deriv}.
The theorems of Section~\ref{sec:sg:thermo} (monotonicity, thermodynamic and zero-temperature limits) apply to the mixed
model in the same way.
\end{proof}

\begin{remark}
Not formalized: the existence of the limit $F(\beta)=\lim_NF_N(\beta)$ (Guerra--Toninelli), which is a hypothesis of
Theorem~\ref{thm:sg:zerotemp}; the Parisi formula; and the almost-sure convergence of the mixed series is shown here
only under the stronger assumption $\sum_p|\beta_p|<\infty$ (implied by Panchenko's $\sum_p2^p\beta_p^2<\infty$,
Lemma~\ref{lem:sg:summable}).
\end{remark}
